% ProbabilityStatistics_A_#9
%%%%%%%%%%%%%%%

%\begin{souprop}
    確率変数$X$の稠密関数は$f$とする。$F\colon \R\to \R$が連続、コンパクトな台を持つとき、
\begin{equation} \label{eq:F:R-> R・連続・コンパクトな台を持つ}
    \mathbb{E}[F(X)]= \int_{-\infty}^\infty F(X) f(x)dx
\end{equation}
%\end{souprop}

%\begin{souprop}
    確率変数$X$の稠密な関数を$f$とする。$F\colon \R\to \R$が連続かつ有界ならば、\eqref{eq:F:R-> R・連続・コンパクトな台を持つ}が成立。
%\end{souprop}
\begin{proof}
    $M=:\sup_{x\in \R} |F(x)|$とかく。$\phi_R\colon \R\to \R$を連続で以下を満たす関数とする。
\begin{equation}
    \phi_R(x) 
    = 
\begin{cases}
    0,\ \ |x|\ge R+ 1\\
    R+ 1- |x|,\ \ R< |x|< R+ 1\\
    1,\ \ |x|\le R
\end{cases}
\end{equation}
\begin{tikzpicture}
    % TikZで図を入れる
\end{tikzpicture}
    このとき$F_R(x)= F(x)\phi_R(x)$はコンパクトな台を持つから、
\begin{equation}
    \mathbb{E}[F_R(X)]= \int_{-\infty}^\infty F(x)\phi_R(x)f(x)dx
\end{equation}
    $F_R(x)\to F(x),\ x\in \R(R\to \infty),\ |F(X)|\ge M$だからLebesgureの優収束定理より
\begin{equation}
    \lim_{n\to \infty}\mathbb{E}[F_n(X)]= \mathbb{E}[F(X)]
\end{equation}
    一方で
\begin{align}
    \left|\int_{-\infty}^\infty F(x)f(x)dx- \int_{-\infty}^\infty F(x)\phi_r(x)f(x)dx\right|
    &= \int_{-\infty}^{-R}|F(x)|(1- \phi_R(x))f(x)dx+ \int_R^\infty |F(x)|(1- \phi_R(x))f(x)dx\\
    &\le M\int_{-\infty}^{-R} f(x)dx+ M\int_R^\infty f(x)dx\\
    &\to 0\ \ (R\to \infty)\ \ \left(\because \int_{-\infty}^\infty f(x)dx= 1\right)
\end{align}
\begin{equation}
    \therefore \lim_{R\to \infty}\int_{-\infty}^\infty F(x)\phi_R(x)f(x)dx= \int_{-\infty}^\infty F(x)f(X)dx
\end{equation}
\end{proof}

%\begin{souprop}
    \ref{prop:Prop1.49}と同じ仮定で$F$は非負連続とする。このとき\eqref{eq:F:R-> R・連続・コンパクトな台を持つ}が成立(有界性は仮定しない)。
%\end{souprop}
\begin{proof}
    $F\colon \R\to \R$は$F(x)\ge 0, x\in \R$を満たす連続関数とする。$n\ge 1$に対し、
\begin{equation}
    F_n(x)
    = 
\begin{cases}
    F(x),\ \ (F(x)\le n)\\
    n,\ \ (F(x)> n)
\end{cases}
\end{equation}
\begin{tikzpicture}
    % TikZで図を入れる
\end{tikzpicture}
    により定める。
    \ref{prop:Prop1.49}から
\begin{equation}
    \mathbb{E}[F_n(X)]= \int_{-\infty}^\infty F_n(x)f(X)dx
\end{equation}
    $F_n(X)$は$n$について単調増加、非負なので
\begin{equation}
    \lim_{n\to \infty} \mathbb{E}[F_n(X)]= \mathbb{E}[F(X)]
\end{equation}
    ここで
\begin{equation}
    \int_{-\infty}^\infty F(x)f(xd)dx< \infty
\end{equation}
    であるとする。$\forall \epsilon> 0$に対して、
\begin{equation}
    \int_{-\infty}^{-R} F(x)f(x)dx+ \int_R^\infty F(x)f(x)dx< \epsilon
\end{equation}
    となる$R> 0$が取れる。また$M_r= \sup_{|x|\le R} |F(x)|$とおき、$n\ge M_R$ならば$F_n(x)= F(x), |x|\le R$
\begin{equation}
    \therefore \int_{-R}^R F(x)d(x)dx= \int_{-R}^R F_n(x)f(x)dx
\end{equation}
    さらに$F_n(x)\le F(x)$なので、
\begin{align}
    \int_{-\infty}^\infty F_n(x)f(x)dx
    &= \int_{-\infty}^{-R} F_n(x)f(x)dx+ \int_{-R}^R F(x)f(x)dx+ \int_R^\infty F_n(x)f(x)dx\\
    &= \int_{-\infty}^\infty F(x)f(x)dx+ \int_{-\infty}^R F_n(x)f(x)fx+ \int_R^\infty F_n(x)f(x)dx- \int_{-\infty}^R F(x)f(x)dx- \int_R^\infty F(x)f(x)dx
\end{align}
\begin{equation}
    \therefore \left|\int_{-\infty}^\infty F_n(x)f(x)fx- \int_{-\infty}^\infty F(x)f(x)dx\right|\le 2\epsilon
\end{equation}
\begin{equation}
    \lim_{n\to \infty} \int_{-\infty}^\infty F_n(x)f(x)dx= \int_{-\infty}^\infty F(n)f(x)dx
\end{equation}
%\begin{souremark}
    Lebesuge積分だと単調収束定理で一発だが，広義積分だと面倒
%\end{souremark}
    また
\begin{equation}
    \int_{-infty}^\infty F(x)f(x)dx= \infty
\end{equation}
    も同様に
\begin{equation}
    \lim_{n\to \infty}\int_{-\infty}^\infty F_n(x)f(x)dx= \infty
\end{equation}
    が示せる。
\end{proof}
%\begin{soucor}
    $F\colon \R\to \R$が連続で$\mathbb{E}[|F(X)|]< \infty$のとき
\begin{equation}
    \mathbb{E}[F(X)]= \int_{-\infty}^\infty F(x)f(x)dx
\end{equation}
    とくに$X$が可積分ならば
\begin{equation}
    \mathbb{E}[X]= \int_{-\infty}^\infty xf(x)dx
\end{equation}
%\end{soucor}
%途中